\section{Conclusion}
\label{sec:conclusion}

This paper studied on-board compression depth and downlink scheduling for
Orbital Edge Computing as one coupled decision problem, and found that the
choice of \emph{quality signal} used to reward that decision matters more
than any single scheduler design choice tested. Grounding utility in
measured downstream segmentation performance rather than pixel fidelity
changed MPC's measured advantage over the best fixed-depth baseline from a
marginal $1.2\%$ to a substantial $32.2\%$, because pixel fidelity is
nearly flat across the compression-depth range studied here while
task-relevant quality is not ($+12\%$ vs.\ $+247\%$). A scheduler is only
as good as the signal it is scored against.

With that signal in place, a single unified utility function reconciled
five previously disagreeing scoring definitions, and a ten-scheduler
comparison under a deliberately routing-sensitive capacity regime showed
that the strongest variant (\texttt{mpc-2level}, a routing MPC and a depth
MPC solved as peers) closes to within $0.5\%$ of the offline optimality
bound. The more consequential finding, however, was architectural rather
than numerical: the $\approx 5$-point utility gap that separates the
strongest and weakest MPC variants tracks \emph{admission control}
(whether a scheduler explicitly drops an overwhelmed task rather than let
it rot in queue), not routing sophistication — the dedicated routing
optimizer's own marginal contribution tops out at $+0.84\%$ utility and
disappears entirely once link rates are not scarce. A related structural
result is that freezing a concrete route for longer than a few minutes
does not survive this constellation's short ground-contact windows,
regardless of how that freeze period is tuned. Finally, a learned (PPO)
policy traded a consistent $5$--$8\%$ utility gap for a
$250$--$400\times$ reduction in per-decision compute time, with an
unanticipated caveat: without a congestion feature in its observation, its
apparent robustness to link scarcity is accidental rather than adaptive.

\paragraph{What remains open.} Several results in this paper are
explicitly flagged, not silently finalized: the uncompressed-reference
mIoU used to anchor the quality signal has not been verified against
FLAIR-1's official held-out test split (\S\ref{sec:downstream}); the scope
of \texttt{mpc-hier} relative to Xuanhao's request to drop certain prior
results is pending his direct confirmation (\S\ref{sec:utility-results});
and both link-rate sensitivity sweeps in \S\ref{sec:sweeps} were generated
against a placeholder quality table and need to be rerun against the real,
measured one before their exact numbers are reported as final. The
proposed satellite-count and ground-station-count sweeps
(\S\ref{sec:sweeps}) and a retraining of the PPO baseline under the
unified reward (\S\ref{sec:rl}) are natural next steps that require no new
code, only additional runs. None of these open items affect the paper's
central conclusion — that a scheduler's measured value depends first on
what it is being asked to optimize, and only second on how cleverly it
optimizes it — but resolving them would sharpen every number built on top
of that conclusion.
